%!TEX program = uplatex
\RequirePackage{plautopatch}
\documentclass[uplatex,dvipdfmx]{beamer}

% テーマとカラーの設定
\usetheme{Madrid}           % スライド全体の世界観（Madrid, Berlin, CambridgeUSなど）
\usecolortheme{default}     % 配色

% 日本語フォント・PDFのしおり文字化け防止
\usepackage{pxjahyper}

% 数学パッケージ
\usepackage{amsmath,amssymb,amsthm}
\usepackage{mathtools}
\usepackage{bm}
\usepackage{graphicx}

% 発表情報の定義
\title{ゼミ・学会発表タイトル}
\subtitle{サブタイトルがあればここに}
\author{あなたのお名前}
\institute{◯◯大学 大学院 理学研究科 数学専攻}
\date{\today}

\begin{document}

% タイトルスライド
\begin{frame}
  \titlepage
\end{frame}

% 目次スライド
\begin{frame}{目次}
  \tableofcontents
\end{frame}

\section{はじめに}

% スライド1枚目
\begin{frame}{研究の背景}
  \begin{itemize}
    \item 箇条書きで背景を説明します。
    \item 1つのスライドに詰め込みすぎないのがポイントです。
    \item 文字サイズを大きめにすると見やすくなります。
  \end{itemize}
\end{frame}

\section{主結果}

% スライド2枚目
\begin{frame}{主定理}
  \begin{theorem}
    任意の $\varepsilon > 0$ に対して、以下が成り立つ……
  \end{theorem}

  \begin{proof}
    証明の概要をここに記述します。
    詳細は配付資料や論文をご参照ください。
  \end{proof}
\end{frame}

\end{document}
